\section{An ordinary involution and four elliptic tests in rank three}
\label{sec:cubic-ordinary-involution}

This section uses the actual definite unitary group and ordinary modules of
J.~Thorne, \emph{A \(p\)-adic approach to the existence of level-raising
congruences}, Proc.\ London Math.\ Soc.\ 128 (2024), e12584.
The precise source locators are (4.1), Section~6, printed pp.~31--36
(the definitions of \(\mathcal A,\mathfrak n_p,\mathfrak t_p,\eta_p,U_p,H_\lambda\),
Lemma~6.1, Proposition~6.2 and Corollary~6.3), Proposition~6.4(2),
and the proof of Theorem~6.11 through (6.3)--(6.4), printed pp.~43--44.
The cohomological inputs remain Thorne's Theorem~5.2 and Theorem~3.5.
The new point is the integral involution at arbitrary continuous weight,
which adds the fourth elliptic test when the rank is three.

Let \(F=F^+F_0\), where \(F_0\) is imaginary quadratic and
\(F/F^+\) is everywhere unramified, and use the group
\[
 G(R)=\{g\in {\rm GL}_n(R\otimes_{\mathcal O_{F^+}}\mathcal O_F):
                  g=(1\otimes c)(g)^{-t}\}.
\]
Assume that \(p\) splits in \(F_0\), and retain all sufficiently-small-level,
split-place, and coefficient-field hypotheses in Thorne's Section~6.
Let \(w_0\) be the antidiagonal permutation matrix.
For a continuous character
\(\lambda:\mathfrak t_p\longrightarrow\mathcal O^\times\), put
\(\lambda^\Theta=w_0(\lambda^{-1})\).
Right translation is \(R_hf(g)=f(gh)\).
The operator on continuous functions is exactly
\[
 U_pf(g)=\sum_{u\in\mathfrak n_p/
                     \eta_p\mathfrak n_p\eta_p^{-1}}f(gu\eta_p),
 \qquad
 \eta_{p,v}=\operatorname{diag}(\varpi_v^{n-1},\ldots,1).
 \tag{G75.1}
\]
There is no weight scalar in (G75.1).
The classical normalized operator is compared with this one by
Proposition~6.2(3); it is not substituted into (G75.1).

\begin{theorem}[AC-G75: integral involution and complete rank-three test]
\label{thm:G75}
The following assertions hold.

\emph{(i) Actual ordinary modules.}
The global involution
\[
 \Theta(g)=w_0(g^{-1})^t w_0=w_0c(g)w_0
 \tag{G75.2}
\]
induces, by \(Cf=f\circ\Theta\), an integral isomorphism
\[
 C:H_\lambda(K^p)\xrightarrow{\ \sim\ }
                H_{\lambda^\Theta}(\Theta K^p)
 \tag{G75.3}
\]
for every continuous \(\lambda\), including nonalgebraic weights.
It transports the away Hecke algebra and its localizations.
At an unramified split place with residue cardinality \(Q\),
\[
 T_i\longmapsto T_n^{-1}T_{n-i}.
 \tag{G75.4}
\]
For the source's Tate-normalized Frobenius roots, this sends the multiset
\(\{\alpha_i\}\) to \(\{Q^{n-1}/\alpha_i\}\).
Thus the residual genericity condition of Theorem~5.2 is preserved.
A prescribed central character, if imposed, is transported to its inverse;
it is not silently held fixed.

\emph{(ii) Four tests.}
Now take \(n=3\), a prime \(q\ne p\) split in \(F_0\), and the distinguished
place \(v_0\mid q\) of Thorne's Section~5.  Retain \(S\cap S_q=\{v_0\}\).
Write \(G_q={\rm GL}_3(F_{\widetilde v_0})\) and
\(G_q^0=\{g:\det g\in\mathcal O_{F_{\widetilde v_0}}^\times\}\).
Let \(\mathfrak m\) be a residual away-Hecke maximal ideal satisfying the
genericity hypothesis of Theorem~5.2.  Fix a locally dominant algebraic
weight \(\mu\) at which this residual system occurs, and suppose there is
no ordinary automorphic representation at that weight and the fixed
away-\(v_0\) level, localized at \(\mathfrak m\), whose \(v_0\)-component is an
unramified twist of \({\rm St}_3\).
For the actual admissible modules
\[
 \mathcal H_\lambda=
 \varinjlim_{V_{v_0}}H_\lambda(K^{p,v_0}V_{v_0})_{\mathfrak m},
\]
one then has, for every continuous \(\lambda\),
\[
 \mathcal H_\lambda[1/p]\otimes^{\mathbf L}_{E[G_q^0]}\pi_\nu=0,
 \qquad
 \nu\in\{(1,1,1),(2,1),(1,2),(3)\}.
 \tag{G75.5}
\]
Here the \(\pi_\nu\) are precisely Thorne's elliptic representations from
Section~3, with trivial central character.

\emph{(iii) An exact fiber suffices.}
Consequently, suppose that for some continuous weight \(\lambda\) there
is a characteristic-zero away-Hecke point \(\mathfrak p\) over
\(\mathfrak m\), occurring in
\(H_\lambda(K^{p,v_0}{\rm Iw}_{v_0})[1/p]\), whose local unramified
Bernstein central character is the cuspidal support of
\(\pi_{1,1,1}\otimes\phi\), for an unramified character \(\phi\).
Then the no-Steinberg hypothesis in (ii) is false.

In particular, the fiber condition holds if the actual determinant of
Proposition~6.4 at \(\mathfrak p\) is unramified at \(\widetilde v_0\) and
has the exact characteristic-zero Frobenius multiset
\[
                  \{a,aQ,aQ^2\},\qquad a\ne0.
 \tag{G75.6}
\]
The equalities in (G75.6) are exact, and are not congruences modulo \(p\).
All field, level, genericity, and actual-Hecke-occurrence hypotheses
above remain part of this implication.
\end{theorem}

\begin{proof}
First, \(w_0\) is a rational unitary matrix and
\(g\mapsto c(g)=(g^{-1})^t\) is an automorphism of \(G\).
Thus (G75.2) is a global automorphism, preserves \(G(F^+)\), and has square
one.  At each chosen split place it is the displayed inverse-transpose
automorphism in the source's matrix coordinates.
It preserves \(\mathfrak n_p\) and \(\mathfrak t_p\).
Put
\[
 z_v=\varpi_v^{-(n-1)}I_n,\qquad z=(z_v)_{v\mid p}.
\]
Then \(z\) is central in \(G(F_p^+)\) and
\(\Theta(\eta_p)=z\eta_p\).
Consequently \(\Theta\) bijects the quotient indexing (G75.1) with itself,
and a direct substitution in that finite sum gives
\[
        C U_p C=R_zU_p,\qquad U_pC=C R_zU_p.
 \tag{G75.7}
\]
These identities hold on the spaces of continuous functions before
taking ordinary parts.

We justify the integral ordinary assertion without assuming that
central translation is a scalar.  Both \(R_z\) and \(R_{z^{-1}}\)
preserve the continuous \(\mathcal O\)-valued functions invariant under
\(G(F^+)\) on the left and \(K^p\mathfrak n_p\) on the right.
They preserve every \(\mathfrak t_p\)-eigenspace and commute with \(U_p\).
If \(P(U_p)f=0\) with \(P\in\mathcal O[X]\) and \(P(0)\) a unit, the
same relation holds for \(R_z^{\pm1}f\).
They therefore preserve the actual ordinary lattice
\(H_\lambda(K^p)\) in both directions.

By Proposition~6.2(1), this lattice is finite free over \(\mathcal O\).
The operator \(U_p\) is an automorphism of this lattice.
One can see this directly from the definition: choose an
\(\mathcal O\)-basis and multiply the finitely many unit-constant
annihilating polynomials for its elements.  The resulting polynomial
annihilates \(U_p\) on the whole lattice and expresses \(U_p^{-1}\) as an
integral polynomial in \(U_p\).
It follows that \(R_zU_p\) is an integral automorphism.
Its characteristic polynomial is monic with unit constant term.
Using the second identity in (G75.7), this polynomial annihilates
\(U_p\) on \(C H_\lambda(K^p)\).
Hence every vector in that image is ordinary in the literal source
sense.

If \(f\) has \(\mathfrak t_p\)-character \(w_0\lambda^{-1}\), then
\[
 (Cf)(gt)
  =\lambda\bigl(w_0\Theta(t)w_0\bigr)^{-1}(Cf)(g)
  =\lambda(t)(Cf)(g).
\]
This is exactly the target character
\(w_0(\lambda^\Theta)^{-1}\).
Thus we have the inclusion in (G75.3), and applying the same argument
again, using \(\Theta^2=1\), proves equality.
The proof applies to every continuous \(\lambda\).

The automorphism sends the level \(K^p\) to its actual image
\(\Theta K^p\).  Sufficient smallness is preserved.
At unramified places it preserves the integral hyperspecial subgroup.
One may equivalently transport the source's fixed hyperspecial
identifications.  At any finitely many exceptional split places the
image level is explicitly used, not identified with the original one.
For a split hyperspecial double coset, inverse-transpose changes
\(\operatorname{diag}(\varpi^i,1^{n-i})\) into the central inverse
uniformizer times the complementary double coset.  This proves
(G75.4).  Substitution into the polynomial
\[
 P_w(X)=\sum_{i=0}^n(-1)^iQ^{i(i-1)/2}T_iX^{n-i}
\]
gives the asserted Frobenius-root transformation.
All ordered ratios are inverted, which preserves the condition that
no ordered ratio equals the cyclotomic character.
This also proves that the transported maximal ideal has the required
genericity.  Since \(C\) intertwines the entire away Hecke algebra,
localization commutes with (G75.3).
For a central element \(h\), \(\Theta(h)=h^{-1}\); this proves the
central-character assertion.

We next prove (ii).  The argument in the proof of Thorne's
Theorem~6.11 before the use of Lemma~6.10 depends only on the stated
genericity, the actual ordinary modules, and the absence of a
Steinberg constituent at the classical weight.
For completeness, its steps are as follows.
Theorem~5.2, followed by the ordinary direct summand, concentrates
\[
 R\Gamma_c(\Omega^{2,{\rm ca}},k)
           \otimes^{\mathbf L}_{k[G_q^0]}\mathcal S
\]
in degree \(2\), where
\(\mathcal S=\mathcal H_\lambda/\varpi\mathcal H_\lambda\) is independent
of \(\lambda\) by Corollary~6.3.
Finite generation and reduction modulo \(\varpi\) imply that the
corresponding integral complex at weight \(\mu\) has one finite free
cohomology group, in degree \(2\).
At this classical weight the characteristic-zero representation is
semisimple and its local constituents are generic, by Corollary~4.3
and its following remark.  Theorem~3.4 then shows that nonzero
degree-\(2\) cohomology would give an unramified Steinberg twist.
The assumed absence of such a constituent therefore makes the integral
complex zero.  Its reduction is zero as well.
Corollary~6.3(3) and the direct-sum decomposition in Theorem~3.5 now give
\[
 \mathcal H_\lambda[1/p]\otimes^{\mathbf L}_{E[G_q^0]}\pi_\nu=0
       \quad(\nu=(1,1,1),(2,1),(3))
 \tag{G75.8}
\]
for every continuous \(\lambda\).
This is exactly the source's (6.3)--(6.4), without invoking its later
two-block local hypothesis.

Apply the same argument to the transported data
\((\Theta K^{p,v_0},\mathfrak m^\Theta,\mu^\Theta)\).
The weight \(\mu^\Theta\) is locally dominant algebraic:
its algebraic part is \(-w_0\mu\).
The source comparison with classical forms is equivariant, and
(G75.3) shows that a Steinberg constituent for the transformed
data would give one for the original data.
Thus the same absence hypothesis holds there.
On \(G_q\), twisting an elliptic representation by \(\Theta\)
reverses its ordered partition.  This follows directly from the
standard-parabolic realization in Section~3 (and agrees with the
smooth-dual formula in the proof of Theorem~3.3).
In particular \((2,1)\) is changed into \((1,2)\), and the other
two partitions in (G75.8) are fixed.
The automorphism preserves \(G_q^0\), so transport of the derived
tensor product is legitimate.
Transporting (G75.8) back by (G75.3) supplies the missing
\((1,2)\) test.  This proves (G75.5).

Finally, suppose the exact fiber in (iii) exists.
We may extend the characteristic-zero coefficient field to an
algebraic closure; the vanishing and nonvanishing assertions in use
are preserved by this flat extension.
Localize the actual smooth module at \(\mathfrak p\), and, if desired,
pass to its residue-field quotient.
The Iwahori invariants of that quotient are nonzero: the invariants
are a finite module over the finite-dimensional away Hecke algebra,
and localization followed by its residue field is nonzero by
Nakayama's lemma.  In characteristic zero, taking Iwahori invariants
is exact.
All the degree-zero tensor vanishings in (G75.5) pass to this quotient
by right exactness.  This is an ordinary quotient argument; no
derived tensor product with a Hecke residue field is assumed exact.

Project to the unramified principal Bernstein component.
It is an actual direct summand, so those degree-zero vanishings
persist.  Call the resulting nonzero module \(M\).
The Iwahori type equivalence identifies \(M\) with its module of
Iwahori invariants.  The latter is finite-dimensional, by the
admissibility in Corollary~6.3.  Hence \(M\) has finite length.
Its irreducible constituents have the prescribed central
character, and are exactly
\(\pi_\nu\otimes\phi\), for the four ordered partitions of \(3\).
Choose an irreducible quotient \(\tau\) of \(M\).
There is therefore a partition \(\nu\) such that
\(\tau^\vee|_{G_q^0}\simeq\pi_\nu|_{G_q^0}\):
unramified twists are trivial on \(G_q^0\), and the four partitions
are closed under contragredients.
The evaluation pairing of an admissible irreducible representation
with its smooth dual is nonzero, so
\[
               \tau\otimes_{E[G_q^0]}\pi_\nu\ne0.
\]
Right exactness for \(M\twoheadrightarrow\tau\) contradicts (G75.5).
Thus the original classical no-Steinberg hypothesis is false.

To justify the stated sufficient Frobenius criterion, use
Proposition~6.4(2), namely the actual equality between the local
group determinant and the symmetric Iwahori-Hecke character.
The exact multiset (G75.6) is, up to a common unramified twist, the
Tate-normalized semisimple parameter of \(\pi_{1,1,1}\).
It therefore gives exactly the Bernstein character used above.
This last step uses no approximation of that polynomial and no
assertion that a residual Hecke point lifts with arbitrarily
prescribed local characteristic-zero roots.
\end{proof}

\begin{remark}[The remaining arithmetic premise]
\label{rem:G75-scope}
The theorem establishes a complete conditional local test on the
actual ordinary modules.  It does not construct the point
\(\mathfrak p\) in part~(iii).
In particular, a residual Frobenius multiset
\((-1,1,-1)\) at \(q\equiv-1\pmod p\) is compatible with (G75.6),
but does not imply that a full-chain characteristic-zero point
occurs in the required ordinary Hecke support.
A classical unramified tempered parameter has equal complex
absolute values after the usual unitary normalization; the ratios
in (G75.6) are different.  Thus the contemplated unramified point
is normally nonclassical, and ordinary interpolation or an actual
support argument is still required.
The theorem also supplies no per-curve containment for the
modified cubic deformation ring.
If an exact fiber is established, the resulting classical
Steinberg occurrence can be fed into the separately proved
projective finite-type transfer, with all of that transfer's
hypotheses retained.
\end{remark}

\begin{remark}[Geometric dictionary and reference scope]
The map \(\Theta\) is an automorphism of the definite unitary
quotient and of its local systems, and (G75.3) is pullback on the
actual integral ordinary lattices.  The four tests are tensor
tests against the elliptic constituents of local
Drinfeld-space cohomology and their transported counterparts.
The only geometric cohomology concentration imported here is the
published Theorem~5.2, with its generic residual hypothesis.
No new assertion is attributed to SGA, and no assertion of
formal verification or of an independently reconstructed
Rapoport--Zink, Caraiani--Scholze, or Dat theorem is made.
The Iwahori type equivalence and classification of the
full-chain principal-series constituents are the standard
characteristic-zero local representation-theoretic inputs also
used in Thorne's Section~3 and proof of Proposition~6.8.
\end{remark}
